
\section{Derivations and differentials}
\section{Separability}
\section{Higher derivations}

\begin{theorem}[Theorem 25.1]\label{mcrt-auto-theorem-25-1}\dcref{25.1}\source{matsumura-commutative-ring-theory:page-0208}\uses{}
\begin{mathematicalrecord}
First fundamental              exact sequence). A composite k L
 A --%B    of ring homomorphisms               leads to an exact sequence of B-
;modules

   (1)     Q,&B-~-?~&RB,A+O.
 where the maps are given by cr(d,,,a 0 b) = bd,,,g(a) and B(d,,,b)= d,,,b
 for aeA and &B. If moreover B is O-smooth over A then the sequence
    (2) 0 -+ R,,, 0 B -         R,,, -   QBiA -+ 0,
 obtained from (1) by adding 0 +at the left, is a split exact sequence.
\end{mathematicalrecord}
\end{theorem}
\begin{proof}
\begin{mathematicalrecord}
In order for a sequence N? AN           -%N? of B-modules to be
?exact, it is sufficient that for every B-module T, the induced sequence
           Hom,(N?,      T) L      Hom,(N,    T) z    Hom,(N?,   T)
is exact. Indeed, taking T = N?, we get a*B*(l T) = 0, and therefore /XX= 0;
and taking T = N/Imcr, we see easily that Ker p = Ima. From this, to
prove that (1) above is exact, it is enough to show that for any B-module
 T,
    (3) Der,(A, T)-     Der,(B, T)+-Der,(B,    T)+O
is exact, but this is obvious.
   Now suppose that B is O-smooth          over A. Choose DEDer,(A,                 T) and
consider the commutative diagram



         gt T
         AABeT
                            with q(a) = (ga, Da).


Then by assumption, there exists k:B -B*T         which can be added to
the diagram, leaving it commutative.     If we write k(b) = (6, D?h) then
D?: B - T is a derivation of B such that D = D?oy, and D? corresponds
to a B-linear map a?:Q,,, -   T. Now take T to be fiZAik@ B, and define
D by D(a) = dAik(a)@ 1, so that D = D?og implies that CL?C(= 1T. Thus (2)
is split. n
   Now consider     the case k A A 5 B when y is surjective; set
Ker g = m, so B = A/m. Then in (1) of the previous theorem we of course
have RsiA = 0, and we want to determine Kercc.
\end{mathematicalrecord}
\end{proof}


\begin{theorem}[Theorem 25.2]\label{mcrt-auto-theorem-25-2}\dcref{25.2}\source{matsumura-commutative-ring-theory:page-0209}\uses{}
\begin{mathematicalrecord}
Second fundamental            exact     sequence).         In   the   above
notation, we have an exact sequence

  (4)   m/m2 LQ       ,wChB     AQ,,d
where 6, is the B-linear map defined by 6(xmodm2)                = dAlkx@ 1. If B is
O-smooth over k then
   (5) O+m/m2-QZ,,,@B-Cl,,,-+0
is a split exact sequence.
\end{mathematicalrecord}
\end{theorem}
\begin{proof}
\begin{mathematicalrecord}
We once more take an arbitrary          B-module     T and consider
  (6) Hom,(m/nt2,     T),   ?    Der,(A,   T)4: il*    Der,(B,     T).

For DEDerdA,     T), to say that 6*(D) = 0 is just to say that D(m) = 0, SO
that D can be considered as a derivation from B = A/m; hence (6) is exact.
If B is O-smooth over k then the extension


of the k-algebra B by m/m2 splits, that is there exists a homomorphism
of k-algebras s:B -+ A/m2 such that gs = 1,. Now sg: A/m2 --P A/m2 is a
homomorphism      vanishing on m/m2, and g(1 - sg) = 0, SO that if we set
A --+ A/m2 5           m/m? 5      T

is an element of Der,(A, T) satisfying S*(D?) = $. Indeed, for xEm, if we
let x = x mod m2 then
         D?(x) = t,b(D(X)) = Ii/(X - sg(X)) = I+@).
Therefore 6* is surjective. If we set T = m/m2 then we see that (5) is a
split exact sequence. n
\end{mathematicalrecord}
\end{proof}


\begin{theorem}[Theorem 25.4]\label{mcrt-auto-theorem-25-4}\dcref{25.4}\source{matsumura-commutative-ring-theory:page-0211}\uses{}
\begin{mathematicalrecord}
Let K be a field of characteristic p, and let 0 # DEDer (K).
    (i) 1,D,D2 ,..., Dpm ? are linearly independent over K;
   (ii) the only way in which c0 + c, D + ... + cp-, Dp- ? with c+K can be
a derivation is if co = cz = ... = cpm 1 = 0.
Pro~$ For UEK, write aL for the operation of multiplying         by a; then
the property D(nx) = D(u).x +a.Dx        of a derivation means that DC-a,=
D(u), + CID. We can write the Leibnitz formula as

         D?oa, = aD?+     i.D(a)D?-?   +

our proof exploits this formula.
    (i) For some i < p suppose that 1, D,. . . , D?- ? are linearly independent
over K, but that l,D ,..., D? are not. Then we can write D? = ci- I Dim ? +
... + co, with c,EK. If we choose some UEK such that D(a) # 0, then in
view of DiouL = ciPl Die?ou, + ..., we get
           aDi + i.D(a)D?-? + ... = ci- ,aD?-?    + ...,
where    . indicates a linear combination of 1, D, . . . , Dim 2. Subtracting     a
times our original relation from this gives a relation of the form
        i.D(a)D?-?   = ... ,
and this contradicts     the assumption      that 1, D, . . . , Die? are linearly
independent.
   (ii) Suppose that E = ciD? + ... + c1 D + c0 is a derivation of K, with i C/J
and ci # 0. Then E(1) = co, so that c0 = 0. Now if i > 1 then take UEK
such that D(u) # 0, and substitute both sides of Eoa, = ciLhclL + ..? in the
Leibnitz formula: we get
         aE + E(a)L = aciDi + [i,q.D(a)      + ac,_,]D?-?    + ..?,
but then in view of the linear independence of 1, D, . . . , Dp- ?, the coefficients
of D?-? on both sides must be equal; therefore i.ci.D(u) = 0, which is a
contradiction.   H

Remarks. (i) The theorem also holds if char K = 0.
   (ii) If K is not a field, this result does not necessarily hold. For example,
let k be a field of characteristic       p, and set A = k[X]/(Xp)   = k[xl, with
$ = 0; then every derivation of k[X] will take the ideal (Xp) into itself,
and therefore induces a derivation of A. In particular, the derivation
Xp-l.a/dX           of k[X]   induces DEDer,(A)       such that D(x)=x~-~,     but
    i) = i.xi-lxP-l       = 0 if i > 1, and therefore  for p > 2 we have 0? = 0.
\end{mathematicalrecord}
\end{theorem}
\begin{proof}
\notready
\end{proof}


\begin{theorem}[Theorem 25.5]\label{mcrt-auto-theorem-25-5}\dcref{25.5}\source{matsumura-commutative-ring-theory:page-0212}\uses{}
\begin{mathematicalrecord}
the Hochschild formula). Let A be a ring of characteristic
p; then for aEA and DED~~(A) we have
         (aD)? = apDP + (aD)?-?(a).D.
\end{mathematicalrecord}
\end{theorem}
\begin{proof}
\begin{mathematicalrecord}
Set E = aD. Then E2 = E~a,~D = (aE + E(a))D = a2 D2 + E(u)D,
and proceeding by induction, we get a relation of the form
                                           k-l
         Ek=ukDk+                           c bk,iDi+Ek-l(u)D,
                                           i=2

where bk,i are elements of A given                               by a pUrdy       fOrma   computation,     so

         bk,i         = fk,i(a,           D(U),   D2(a), . . .) D?-?(a)),
where the fk,i are polynomials with coefficients in Z/(p) not depending on
A, on a or on D. Now to prove our theorem we need only show that ,j?,,i = 0
for 1 < i < p. Let k be a field of characteristic p, and let x1, x2,. . . be a
countable number of indeterminates over k; set K = k(x,, x2,. .). Define
a k-derivation D of K by Dxi = xi+ r. (Since C12,,,is the free K-module with
basis dx,, dx,,. .., given any ,fieK there exists a unique DEDer,(K)
such that Dxi = fi.) For this D, we set E = x,D; then since EP- xpDp =
bp,p-lDP-l + ... + b,,2D2 + EP-l(u).D is a derivation, by the previous
theorem we must have b,,i = 0 for 1 < i < p. Therefore
         bp,i=fg,i(X1,X2,...,Xp~i+1)=0,


and this proves that f,,i = 0. w
  This formula is known as the Hochschild formula, although it is also
reported to have been first proved by Serre. Be that as it may, it is an
important fact that (uD)? is a linear combination of Dp and D.
\end{mathematicalrecord}
\end{proof}


\begin{theorem}[Theorem 26.3]\label{mcrt-auto-theorem-26-3}\dcref{26.3}\source{matsumura-commutative-ring-theory:page-0215}\uses{}
\begin{mathematicalrecord}
If k is a perfect field then every extension field K of k is
 separable over k, and a k-algebra A is separable if and only if it is reduced.
\end{mathematicalrecord}
\end{theorem}
\begin{proof}
\begin{mathematicalrecord}
Recall that a field k is perfect if every algebraic extension of k is
 separable. Ifk has characteristic 0 then every extension field K is separably
 generated, and therefore separable. In characteristic p, perfect implies
 k = kl?p, so that by the previous theorem, all sublields of K finitely
generated over k are separable, so that K itself is separable over k. (Caution:
K may fail to be separably generated over k; for a counter-example, let x
be an indeterminate over k, and K = k(x, xp-l, xp-?, . .).) Now we show
that if A is a reduced k-algebra then A is separable. We can assume that
A is finitely generated over k. Then A is a Noetherian ring, and the total
ring of fractions K of A is of the form K = K, x ... x K, by Ex. 6.5. Each
Ki is separable over k, so that K is also separable, and hence so is its
 subring A. n
   In general two subfields K, K? of a given field L are said to be linearly
disjoint over a common subfield k if the following three equivalent
conditions are satisfied:
   (a)ifa,,...,     CI,EK are linearly independent over k they are also linearly
independent over K?;
   (b) the same with K and K? interchanged;
   (c) if we write K [K?] for the subring of L generated by K and K?j the
natural map K OkK? -+ K[K?] is an isomorphism.
Proofofequivalence.        (a)*(c) Let 4 = Cyxi o yi be an element of the kernel
of K OkK? - K[K?].             Suppose that x1,. . , x, are linearly independent
over k, and that the remainder x,, l,. . . , x, arc linear combinations of them7
and rewrite 4 = C; xi myi. The image in K[K?] of 5 is CxiY:, but if this is ?
then by (a) we have y; = 0 for all i, so that 4 = 0. This proves (C).
   (c)+(a)      is also easy; finally, since (c) is symmetric in K and K 12 we Of
course also get (a+(b).
~,,EK are linearly independent over k. If
                               we set k, = k(cl,. . . , {,), so that k, is a finite
 xtension of k; for some sufficiently large n we have kp? c k, and if we set
   = Kgkk,,      then A is a reduced ring. However, A is finite as a K-module,
    is a zero-dimensional ring, but the p?th power of any element of A is in K,
     that we see that A has only one prime ideal. Hence A is a field, and
   N K[k,]. From this we get x&i 0 ci = 0, that is ti = 0 for all i.
   (ii) If K and kP?? are linearly disjoint, then kPm? c kPm?,and hence K and kP-?
are also linearly disjoint over k, so that K QkkPm? is a field. If K? is a subfield
of K which is finitely generated over k then by Theorem 2, K? is separable
over k. Hence K is also separable over k. w



Let K be an extension field of a field k; then QKjk is a vector space over
K, and is generated by {dx(xeK},              so that there exists a subset B c K
such that {dxl xEB} forms a basis of the vector space Q,,,. A subset
Bc K with this property is called a differential basis of K over k. The
following condition (*) is necessary and sufficient for a subset {x~}~~~ c K
to form a differential basis for K over k.
   (*) if ~,EK are specified for every LEA in an arbitrary way, then there
exists a unique DEDer,(K) such that D(x,) = y, for all i.
    For x~,...,x,EK,       let us study the condition for dxl,...,dx,ERKik      to be
linearly independent over k. If k has characteristic 0 then this is equivalent
            x, being algebraically independent over k. Indeed, if there exists
                  1,. . . ,X,] such that J-(x,, . . . ,x,) = 0, then choose such a
relation of smallest degree; suppose for instance that X, actually appears in
f, SO that f1 = af/a X, is non-zero, but of smaller degree than f, and hence
f 1(x) # 0. Then f(x) = 0 gives
           0 = df = 1 fi(x)dxi,
                    ,dx, are linearly dependent. Conversely, if x1,. . . ,x, are
algebraically independent over k then there exists a transcendence basis
B of K/k containing        these, so that there exists k-derivations Di of the Purely
transcendental        extension k(B) satisfying
              Di(Xi) = 1 and Of(v) = 0 for Xi #DEB,
(namely, a/ax,). Moreover, K is a separable algebraic extension of k(D),
and so by Theorem 25.3 is 0-etale, so that the derivations Di extend to
derivations from K to K. Then since Di(xj) = 6ij, the differentials
dx 1, . . . , dw%,k       are clearly linearly independent. Thus in this case a
differential basis is the same thing as a transcendence basis.
    Now consider a field k of characteristic p. We say that elements
X1,...,-   Y,EK of an extension field K are p-independent over k if
 CKW, xl,. , x,):KP(k)]          = p?, and a subset B c K is p-independent if any
finite subset of B is p-independent. This condition means precisely that the
set
                               x1,. . . ,x, are distinct
                               elements of B and 0 6 Cli < p
is linearly independent over F?(k); the elements of rB are called the
p-monomials of B. If B is not p-independent, we say it is p-dependent.
The condition of p-independence is not just a property of B and k, but
also depends on K. If B c K is p-independent over k and K = KP(k,B),
we say that B is a p-basis of K/k. If C c K is p-independent over k then
one can easily show by Zorn?s lemma that there exists a p-basis of K/k
 containing C.
    One sees easily that B is a p-basis of K/k is equivalent to IB being a
basis of K over K?(k) in the sense of linear algebra. If this holds then any
map D: B-K             has a unique extension to an element DEDer,(K). Indeed,
for a p-monomial of B we set
           D(x?,?. . . x~~)   = Cr=   1 OriX9?.   . . X~I-   ?. . . x~?D(x~),

and extend D to K as a KP(k)-linear map; then D is a k-derivation. Thus
a p-basis B is a differential basis of K/k. Conversely, if B? is a differential
basis of K/k then B? is p-independent over k; for if xi,. ., x,EB? are
p-dependent, we can assume that xl~KP(k, x2,. ., x,), so that we can
write x1 = S(x,, . . . , x,), where f is a polynomial with coeflicients in KP(k).
Then in QKjk we get dx, = CZ(af/axi)d x i, which contradicts the linear
independence of dx, , . . , dx,. Now if we take a p-basis B of KJk containing
B?, then since both B and B? are differential bases, we have B = B?. We
summarise the above as follows:
\end{mathematicalrecord}
\end{proof}


\begin{theorem}[Theorem 26.5]\label{mcrt-auto-theorem-26-5}\dcref{26.5}\source{matsumura-commutative-ring-theory:page-0217}\uses{}
\begin{mathematicalrecord}
The notion ofdifferential basis coincides with transcendence
basis in characteristic 0, and with p-basis in characteristic p.
  Now we look at the relation between separability and differential bases.
Letting II c k denote the prime subfield of k, we write Qk for !&,.
\end{mathematicalrecord}
\end{theorem}
\begin{proof}
\notready
\end{proof}


\begin{theorem}[Theorem 26.8]\label{mcrt-auto-theorem-26-8}\dcref{26.8}\source{matsumura-commutative-ring-theory:page-0219}\uses{}
\begin{mathematicalrecord}
Let K/k be a separable extension of fields of characteristic
p, and let B be a p-basis of K/k. Then B is algebraically independent over
k, and K is 0-etale over k(B).
\end{mathematicalrecord}
\end{theorem}
\begin{proof}
\begin{mathematicalrecord}
Suppose by contradiction       that b,, . . . , ~,EB are algebraically
dependent over k. Suppose that 0 # fEk[X,,       . . . ,X,1 is a polynomial of
minimal degree among all those with f(b) = 0, and set deg f = d. Then write
         f(xl~~~~~Xn)=o~i~i                    ipgil..,i,(XE,...rX~)X;I...X~

                                 .    ,r..rn


Then since b,, . . . , b, are p-independent over k and f(b) = 0, we have
gi ,,.,i,(bP) = 0 for all values of i,, . . , i,. However, since
         d 3 deggi,...i,(XP)         + il + .?. + i,,
by choice of f we must have f(X) = g,,,e(XJ?). Hence we can write f in
the form f(X) = h(X)P, with hEkl?p[Xl,.      . ,XJ. However, since K and
klip are linearly disjoint over k, the monomials of degree < d in bl , . . . , bn,
being linearly independent over k, must also be linearly independent over
kllp. Thus h(b) # 0, but this contradicts h(b)P = f(b) = 0. For the second
half, see the proof of the following theorem.      n
\end{mathematicalrecord}
\end{proof}


\begin{theorem}[Theorem 26.9]\label{mcrt-auto-theorem-26-9}\dcref{26.9}\source{matsumura-commutative-ring-theory:page-0219}\uses{mcrt-auto-theorem-25-1}
\begin{mathematicalrecord}
If K is a separable field extension of a field k, then K is
O-smooth, and conversely.
\end{mathematicalrecord}
\end{theorem}
\begin{proof}
\begin{mathematicalrecord}
Let B be a differential basis of K/k. If K/k is separable then by
Theorem 5 and Theorem 8, k(B) is purely transcendental Over k. Then?
as one sees at once from the definition, k(B) is O-smooth over k. Moreover7
KIlJR\ is O-etale. In characteristic 0 this follows from Theorem 25.3; in
j   ?26     Separability                                                                          205

  characteristic p, by Theorem 6 we have an exact sequence O-tR,O K
   -aK       -+fiZKlk+O,  so that putting together an absolute p-basis of k
; with B we have an absolute p-basis of K, and this is clearly also an absolute
  p-basis of k(B). Hence K/k(B) is 0-etale by Theorem 7. Therefore K/k is
; O-smooth.
     Conversely, if K/k is O-smooth then by Theorem 25.1, QR,@ K -0,
  is injective, so that by Theorem 6, K/k is separable. n
5

$ Imperfection    modules and the Cartier           equality
FQtg ui e enerally, if k --+ A + B are ring homomorphisms,      we write rsiAik
i for the kernel of nAikO,,,B -fiZejk,    and call it the imperfection module
1; of the A-algebra B over k. If k = Z or k = Z/p (the prime field of
1 characteristic p) we omit k, and write IBIA.
1
1: Lemma 1. Let k + K -L         --+ L? be field homomorphisms.      Then there
1, exists a natural exact sequence
.*
            0 + rLiKlk 0, c - rLsIKIk - hk
               - QLjrc &Ii -t&f,,      - QL?,L + 0.

    Proof We have a commutative            diagram with exact rows.
            04-L,K,kc3Lfi -qlkcw                    -~L,k~L~            -Q,,,cW+o
                     II 1                      II              12I             f3 1
            O-+      r   L?/K/k     -%,k   @ L?        -         %?,k      -     clL.,I;   -+o.

    We abbreviate this as
            0+X-A-B-P-,0
                   i   II             I    1
            O+Y-A-C-Q+O.
    and from it construct
            o?Ai?-ia+o

            O-A/Y        -C-Q-to;
    applying the snake lemma gives the exact sequence
              O-+Y/X+Kerf,+Kerf,+O,
    from which we easily get
              0+X--+    Y-Kerf,     +P+Q-+Cokerf,+O.
    This is just what we wanted to prove. n
\end{mathematicalrecord}
\end{proof}


\begin{theorem}[Theorem 26.10]\label{mcrt-auto-theorem-26-10}\dcref{26.10}\source{matsumura-commutative-ring-theory:page-0220}\uses{}
\begin{mathematicalrecord}
the Cartier equality). Let k be a perfect field, K an
    extension of k and I. a finitely generated extension field of K; then
       (*I   rb.QL,K = tr.deg,L + rkLrLIKIk,
    (where rk, denotes the dimension of a vector space over L).
\end{mathematicalrecord}
\end{theorem}
\begin{proof}
\begin{mathematicalrecord}
Suppose that k c K c L c L?, with L finitely generated over K and
L? finitely generated over L. If the theorem holds for k c L c L? and for
kc K c L then both rEILlk and r,,K, 0 L are finite-dimensional    over L,
so that by the lemma, f-L,,K,L is also finite-dimensional, and
           rkL,QLs,K - rkL,,K,k
              = (rk,,&,,, - rk,G~,L,k) + (rk,%,, - rWLiKik)
              = tr.deg,L?+    tr.deg,L= tr.deg,L?;
thus the theorem also holds for k c K c L?. Now every finitely generated
field extension can be obtained by a succession of the following three
kinds of simple extensions:
    (1) L = K(a) where c1is transcendental over K;
    (2) L = K(E) where c1is separable algebraic over K;
    (3) L = K(cc) where char K = p, and G??= aeK, but cr$K.
Hence we need only prove (*) in each of these special cases. (1) and (2)
are easy. For (3), if we write L = K[X]/(Xr        - a) we see that
           fhlk = (QKLXIIk 0 WLda
                = (R&K     da) 63 L 0 Lda,
and dcc# 0. Furthermore, since k is a perfect field, we have a$ Kr = kKP,
so that in QZKlk(= 0,) we have da # 0, rk R,,, = 1 and rk rLIKIL = 1, so that
(*) also holds in this case. n
\end{mathematicalrecord}
\end{proof}


\begin{theorem}[Theorem 27.1]\label{mcrt-auto-theorem-27-1}\dcref{27.1}\source{matsumura-commutative-ring-theory:page-0223}\uses{}
\begin{mathematicalrecord}
If the ring A is O-smooth over a ring k, then a higher
derivation of length m < cc over k from A to an A-algebra B can be
extended to a derivation of length GO.

  This theorem can be applied for example to the case of a field k and a
separable extension field A of k.
If A is a ring of characteristic p then an ordinary derivation of A is
  zero on the subring AP, but higher derivations need not vanish on AP. If
  Q = (Do, D,, D,, . .) is a higher derivation of length m > p then from
 E,(d) = E,(a)P = up + D1(a)P.tP + ..., we get
            Dp(aP) = D, MP, and in general D,,(&) = D, (u)p?.
 For example, it follows from this that if k is a field of characteristic p, and
 K= k(a) with aPEk but a$k, then although there exists DEDerk(K)
 such that D(a) = 1, this D cannot be extended to a higher k-derivation of
 length B p. Since K is separable over the prime sublield, D can be extended
 to a higher derivation of length GO (over the prime subfield), but this
 extension cannot be trivial on k.
    We say that a higher derivation g = (Do, D,, . . .)EHS,(A) is iterative if
 it satisfies the following conditions:

                           Di+ j   for all    i,j.

 This condition   is equivalent    to asking that the following     diagram    is
 commutative:




  where E,(u) = C t?D,( a) and E,(x tYuv) = C t?E,(u,), and the right-hand
: vertical arrow is the inclusion map. Indeed,
           U-W))     = E,(~~?D&))    = 1 t?x u?D,D,W>
                                        ? B
: and

            E,+,(u) = T(t + u)?D,(u) = 1 t?x u?                 D,+,(u).
                                               v fl
;. If A contains the rational field Q, then one sees by induction on n that
   an iterative higher derivation            satisfies D, = D:/n!,   and is hence
  determined by D, only. Conversely, for DEDerk(A), we see that
  (1, D, D2/2!, D3/3!, . . .) is an iterative higher derivation. If A has charac-
i teristic p then for an iterative D = (D,, D,, . . .) we have Di = Di/i! for
; i <P, and DI = 0. Thus one cannot hope to extend any derivation to an
1 iterative higher derivation, even if A is a field.

i Th eorem 27.2. Let k LA 2 B be ring homomorphisms,    and suppose
? that B is 0-etale over A. Then given D = (D,,D,, . . .)EH&(A, B),
 there exists a unique D? = (Db,D;,          . .)EH?~(B)   such that Di(g(a)) =
D,(a) for all i. Moreover,    if D* = (D;F,DT,. ..)eHS,(A)      is such that
Di = goD* for all i, and if D* is iterative, then D? is also iterative.
\end{mathematicalrecord}
\end{theorem}
\begin{proof}
\begin{mathematicalrecord}
There is no problem about De = 1,. Now assume that (Db,. ..,D&)E
HS,(B, m) has been constructed so that Diog = Di for i < m; then if we define
                                                                           m+l
         h:A   -B,+            1   = B[t]/(t??)               by h(a) = c        t?DJa),

andu:B--+B,byu(b)=                    C~t?D~,(b),             we obtain the left-hand commutative
diagram.




                                                          h
          A h--?&+1                               A   -          B,+,


Hence by the 0-etale assumption, there exists a unique o:B -+ B,, , which
makes the right-hand diagram commutative. Repeating this we see that
D exists and is unique. If r>* is iterative, and we consider the homo-
morphisms    E,: A -+ A[r] and E;:B - B[t] corresponding      respectively
to Q* and D?, then we know E,o E, = E,,,, and we need only prove that
E;oE; = Ej,,. By induction on m, assume that
         EL(E;(b)) = E:+,,(b)mod(t,u)?+?                          for all bEB;
then from the commutativity                  of
                            C.?
          B           -                           But, uj/(t, u)?+ ?

          I,                                        II
          A -f%           A[t,uj      --+B[t,u]/k4m+2,

from E,,, = E,oE, and from the assumption                                that B is 0-etale over A, we get
         EI(E;(b)) = E;+,(b) mod (t, u)~+~                              for all beB.
This proves that EuoE; = E;,,.    m
\end{mathematicalrecord}
\end{proof}


\begin{theorem}[Theorem 27.3]\label{mcrt-auto-theorem-27-3}\dcref{27.3}\source{matsumura-commutative-ring-theory:page-0225}\uses{}
\begin{mathematicalrecord}
(i) Let A be a ring of characteristic            p, and suppose
that XEA, DEDer(A)           satisfy Dx= 1 and Dp=O;           set A, = {~EAI
Da = 0). Then A is a free module over A,, with basis 1, x,. . , xp- ?.
   (ii) Let k be a field of characteristic p, and K a separable extension of k;
let DEDer,(K)      be such that DfO,           DP=O. Set K,={a~KlDa=o).
Then there exists an XE K such that Dx = 1, and a subset B, c K, such
that B = {x} uB, is a p-basis for K over k.
\end{mathematicalrecord}
\end{theorem}
\begin{proof}
\begin{mathematicalrecord}
(i) Suppose that a0 + zlx + ... + xixi = 0 for some i < p and ccje,4,.
Then applying D? we get i!a, = 0, hence ri = 0. Hence by downward
induction on i we see that 1, x, . . . , xp- ? are linearly independent over A,.
Now in view of Dp = 0, for every aeA we have D?+ ?a = 0 for some
0 < i <p. If i = 0 we have LIEA,,. If i > 0 then D?(a - x?D?a/i!) = 0, so that we
see by induction that
         D??u   = O=aEA,      + A,x + ... + A,x?;
?Setting i = p - 1 in this gives A = A, + A,x + ... + A,xP-?.
    (ii) Since D # 0 we can find ZCSK such that Dz # 0. Now in view of
DPz = 0, there is an i such that D?z # 0 but D?+?z = 0. If we set y = D?z
 and x = (D?- ?z)/y then Dx = 1, so that by (i) we have K = K,(x) and
[K:K,J =p. Now if we had xPgKP,k, then x~K,k?~,                    and we could
write x=x;wiai        with q,...,       w,EK~ linearly independent over k, and
aEk?lP. Now kc K, and x#K,,                   so that x,ol ,..., o, are linearly
independent over k, and hence by the assumption that K is separable over
k, they are also linearly independent over kllp, which contradicts x = 1 qq.
 Thus xP$Kgk. Hence we can choose a p-basis C of K, over k such that
 x~EC; set B0 = C - {xp). Then if y,, . . . ,y, are distinct elements of B,,
 we have [KP,k(xp,y,,...,y,):K$k]          = p?+l, and together with K = K,(x)
 this gives [KP k(y , , . . . , y,):KPk] =p?. Thus B, as a subset of K is
p-independent      over k. Since also K, = KP,k(xP, B,), we have K =
K,(x) = KPk(x, II,), so that setting B = B, u {x} we get a p-basis of K
 over k. n
\end{mathematicalrecord}
\end{proof}


\begin{theorem}[Theorem 27.4]\label{mcrt-auto-theorem-27-4}\dcref{27.4}\source{matsumura-commutative-ring-theory:page-0226}\uses{}
\begin{mathematicalrecord}
Let K be a field of characteristic p, and k c K a subfield
such that K is separable over k. Then a necessary and sufficient condition
for DEDer,(K)    to extend to an iterative element of H&(K) is that Dp = 0.
\end{mathematicalrecord}
\end{theorem}
\begin{proof}
\begin{mathematicalrecord}
We have already seen necessity, and we prove sufficiency. We can
assume that D #O; if DP = 0 then we can choose K,,x and B, as in
Theorem 3, (ii). We set K? = k(B,); then D is a K?-derivation,    K is 0-etale
over K?(x), and K?(x) is a purely transcendental extension of K?. We define
a homomorphism       @K?(X) -K?(x)[t]      by setting E,(a) = a for CCEK? and
J%(X) = x + t; then
          E?(E,(X)) = x + 24+ t = E,+,(x),
so that E,oEt= Et+u holds over the whole of K?(x). Thus E, defines a
iterative higher derivation D of K?(x) over K?. Since K is O-&ale over K?(x),
bY Theorem 2 there is an extension of D to an iterative higher derivation of
K over K?; the term of degree 1 in D is D, so that Q is an extension of D (or
more precisely, of (1, D)). n
\end{mathematicalrecord}
\end{proof}
